% 第5章 放大电路的反馈
\chapter{放大电路的反馈}

反馈是模拟电路中最重要的结构思想：它把输出量引回输入端，用"比较"产生的净输入量去控制输出。本章的逻辑是：先判断有无反馈、极性、交直流、组态（四个判断），再给出方块图与一般表达式，然后在深度负反馈条件下把闭环增益化为 $1/F$，最后定量讨论负反馈对性能的影响。

\section{反馈的基本概念}

\paragraph{定义}
在电子电路中，将输出量（输出电压或电流）的一部分或全部通过一定的电路形式作用到输入回路，用来影响其输入量的措施称为\textbf{反馈}。

\paragraph{基本组成}
反馈放大电路分为\textbf{基本放大电路}与\textbf{反馈网络}两部分：前者放大信号，后者传输反馈信号。基本放大电路的输入信号称为\textbf{净输入量}，它由输入量与反馈量共同决定，这是理解全部反馈问题的关键。

\paragraph{正反馈与负反馈}
\begin{itemize}
    \item 使基本放大电路\textbf{净输入量增大}的反馈称为\textbf{正反馈}；使净输入量减小的称为\textbf{负反馈}；
    \item 也可以从输出量的变化来判断：反馈使输出量的变化增大为正反馈，使输出量变化减小为负反馈。
\end{itemize}
负反馈用于改善放大电路性能，正反馈用于产生振荡（第 6 章）。

\paragraph{瞬时极性法}
判断极性的标准方法：
\begin{enumerate}
    \item 假定输入信号在某一时刻对地的极性（如"+"）；
    \item 沿信号传输路径逐级推出各点电位、电流的极性，得到输出信号的极性；
    \item 由输出端经反馈网络推回输入端，得到反馈信号的极性；
    \item 比较反馈信号与输入信号在输入端的叠加结果：使净输入量减小的为负反馈，增大的为正反馈。
\end{enumerate}
注意"比较"发生在输入端：串联反馈比较的是\textbf{电压}（同一点电位相加），并联反馈比较的是\textbf{电流}（同一节点电流相加），二者的符号规则不同。

\paragraph{直流反馈与交流反馈}
只在直流通路中存在的反馈称为\textbf{直流反馈}，只在交流通路中存在的反馈称为\textbf{交流反馈}；也可按反馈量的成分判断（只含直流量则为直流反馈）。判断方法是看反馈通路中是否有隔直电容（或旁路电容）：
\begin{itemize}
    \item \textbf{直流负反馈}的作用是\textbf{稳定静态工作点}；
    \item \textbf{交流负反馈}的作用是\textbf{改善动态性能}（稳定增益、展宽频带、减小失真、改变输入输出电阻）；
    \item 二者常常同时存在（如发射极电阻 $R_e$ 未并旁路电容时，既稳 $Q$ 点又引入交流负反馈）。
\end{itemize}

\paragraph{局部反馈与级间反馈}
只跨接在某一级内部的反馈为\textbf{局部反馈}；跨接在两级或多级之间的为\textbf{级间反馈}。分析时先看级间反馈（它决定整体性能），再看局部反馈。

\section{负反馈放大电路的四种组态}

\subsection{四种组态的判断}

组态由\textbf{取样方式}与\textbf{比较方式}组合而成，共四种：电压串联、电流串联、电压并联、电流并联。

\paragraph{取样方式：电压反馈还是电流反馈（输出短路法）}
令输出电压为零（把输出端短路到地），观察反馈量是否消失：
\begin{itemize}
    \item 反馈量随输出电压消失而消失，为\textbf{电压反馈}（反馈信号与输出电压成正比）；
    \item 反馈量仍然存在，为\textbf{电流反馈}（反馈信号与输出电流成正比）。
\end{itemize}

\paragraph{比较方式：串联反馈还是并联反馈}
看反馈信号以何种方式加到输入回路：
\begin{itemize}
    \item 反馈信号以\textbf{电压}形式与输入电压求差得到净输入电压，为\textbf{串联反馈}（反馈网络与输入回路串联，接在不同节点）；
    \item 反馈信号以\textbf{电流}形式与输入电流求差得到净输入电流，为\textbf{并联反馈}（反馈网络与输入回路并联，接在同一节点）。
\end{itemize}
工程判据：反馈支路接到输入管的\textbf{发射极}（或源极）一般是串联反馈，接到\textbf{基极}（或栅极）一般是并联反馈。

\subsection{四种组态的特点}

\begin{table}[htbp]
\centering
\caption{四种负反馈组态及其作用}
\begin{tabular}{@{}lllll@{}}
\hline
组态 & 取样量 & 比较量 & 稳定量 & $R_i$ / $R_o$ \\
\hline
电压串联 & $u_o$ & 电压 & 输出电压 & 增大 / 减小 \\
电压并联 & $u_o$ & 电流 & 输出电压 & 减小 / 减小 \\
电流串联 & $i_o$ & 电压 & 输出电流 & 增大 / 增大 \\
电流并联 & $i_o$ & 电流 & 输出电流 & 减小 / 增大 \\
\hline
\end{tabular}
\end{table}

记忆规律：\textbf{串联增大输入电阻，并联减小输入电阻；电压减小输出电阻，电流增大输出电阻。}稳定的是"被取样的那个量"：电压反馈稳定输出电压，电流反馈稳定输出电流。

\paragraph{易错点}
\begin{itemize}
    \item 电压反馈与电流反馈看的是\textbf{取样点}，不是"输出端有没有负载电阻"；
    \item 判断串并联不能只看"输入信号加在哪个电极"，还要看反馈信号回到哪个电极，两者共同决定比较方式；
    \item 四种组态名称中"电压/电流"在前、"串联/并联"在后，与"取样/比较"一一对应，不能颠倒。
\end{itemize}

\subsection{方块图与一般表达式}

把反馈放大电路抽象为三个环节：基本放大电路（增益 $\dot A$）、反馈网络（反馈系数 $\dot F$）、比较环节（求差）。
\begin{equation}
\dot A=\frac{\dot X_o}{\dot X_i'},\qquad
\dot F=\frac{\dot X_f}{\dot X_o},\qquad
\dot X_i'=\dot X_i-\dot X_f
\formulaNote{$\dot A\dot F$ 为环路放大倍数}
\end{equation}
消去 $\dot X_i'$、$\dot X_f$，得闭环放大倍数
\begin{equation}
\dot A_f=\frac{\dot X_o}{\dot X_i}
=\frac{\dot A}{1+\dot A\dot F}
\formulaNote{$1+\dot A\dot F$ 为反馈深度}
\end{equation}
在中频段各量均为实数，可写为 $A_f=A/(1+AF)$。由该式可得三条基本结论：
\begin{itemize}
    \item $|1+AF|>1$：$|A_f|<|A|$，为\textbf{负反馈}；
    \item $|1+AF|<1$：$|A_f|>|A|$，为\textbf{正反馈}；
    \item $1+AF=0$：$A_f\to\infty$，电路无输入也有输出，即\textbf{自激振荡}。
\end{itemize}

\section{深度负反馈放大电路的分析}

\subsection{深度负反馈的条件与实质}

当 $|1+AF|\gg1$ 时称为\textbf{深度负反馈}，此时
\begin{equation}
A_f=\frac{A}{1+AF}\approx\frac{1}{F}
\formulaNote{闭环增益几乎只取决于反馈网络}
\end{equation}
其物理实质是\textbf{净输入量近似为零}：
\begin{itemize}
    \item \textbf{串联负反馈}：净输入电压 $u_{id}=u_i-u_f\approx0$，故 $u_i\approx u_f$（称为"虚短"）；
    \item \textbf{并联负反馈}：净输入电流 $i_{id}=i_i-i_f\approx0$，故 $i_i\approx i_f$（称为"虚断"）。
\end{itemize}
深度负反馈下闭环增益由反馈网络决定，因此增益的稳定性、精度都取决于反馈元件，这就是"用电阻精度换增益精度"。

\subsection{四种组态下的闭环增益}

四种组态的差别只在于 $F$ 的物理量不同，从而 $A_f$ 的量纲不同：

\begin{table}[htbp]
\centering
\caption{深度负反馈下四种组态的反馈系数与闭环增益}
\begin{tabular}{@{}llll@{}}
\hline
组态 & 反馈系数 & 闭环增益 & 量纲 \\
\hline
电压串联 & $F_{uu}=u_f/u_o$ & $A_{uf}\approx1/F_{uu}$ & 电压放大倍数 \\
电压并联 & $F_{iu}=i_f/u_o$ & $A_{rf}\approx1/F_{iu}$ & 互阻（$\Omega$） \\
电流串联 & $F_{ui}=u_f/i_o$ & $A_{gf}\approx1/F_{ui}$ & 互导（S） \\
电流并联 & $F_{ii}=i_f/i_o$ & $A_{if}\approx1/F_{ii}$ & 电流放大倍数 \\
\hline
\end{tabular}
\end{table}

求电压放大倍数时，若所得闭环增益不是电压比，还需借负载关系换算（如 $u_o=i_oR_L$）。步骤固定为：\textbf{判组态 $\to$ 写出 $F$ 的定义式 $\to$ 用虚短/虚断求 $F$ $\to$ 取倒数得 $A_f$ $\to$ 按需换算}。

\subsection{理想集成运放的性能指标}

把集成运放理想化，可大大简化深度负反馈电路的分析：
\begin{itemize}
    \item 开环差模增益 $A_{od}\to\infty$；
    \item 差模输入电阻 $r_{id}\to\infty$，输出电阻 $r_o\to0$；
    \item 共模抑制比 $K_{CMR}\to\infty$，上限截止频率 $f_H\to\infty$（无频率失真）；
    \item 输入失调电压 $U_{IO}$、输入失调电流 $I_{IO}$ 及其温漂均为零。
\end{itemize}
由 $A_{od}\to\infty$ 而 $u_o$ 为有限值，得 $u_P-u_N=u_{id}=u_o/A_{od}\to0$，即\textbf{虚短}；由 $r_{id}\to\infty$ 得 $i_P=i_N=0$，即\textbf{虚断}。"虚短、虚断"是分析所有理想运放电路的两把钥匙，它们只在运放工作于线性区（有负反馈）时成立。

\section{负反馈对放大电路性能的影响}

\subsection{稳定放大倍数}

对 $A_f=A/(1+AF)$ 求微分，得
\begin{equation}
\mathrm{d}A_f=\frac{\mathrm{d}A}{(1+AF)^2}
\end{equation}
于是闭环增益的相对变化量与开环的相对变化量之比为
\begin{equation}
\frac{\mathrm{d}A_f}{A_f}=\frac{1}{1+AF}\cdot\frac{\mathrm{d}A}{A}
\formulaNote{减小到开环时的 $1/(1+AF)$}
\end{equation}
即反馈深度越大，增益越稳定。代价是增益数值下降为原来的 $1/(1+AF)$，这是"以增益换稳定"。

\subsection{展宽频带}

设高频段开环增益为 $\dot A_h=A_m/(1+\mathrm{j}f/f_H)$，引入负反馈后
\begin{equation}
\dot A_{hf}=\frac{\dot A_h}{1+\dot A_h\dot F}
=\frac{A_{mf}}{1+\mathrm{j}\dfrac{f}{f_{Hf}}}
\end{equation}
\begin{equation}
f_{Hf}=(1+A_mF)f_H
\formulaNote{下限频率降低为 $f_{Lf}=f_L/(1+AF)$}
\end{equation}
因此通频带被展宽：
\begin{equation}
f_{bwf}=f_{Hf}-f_{Lf}\approx(1+AF)f_{bw}
\formulaNote{带宽展宽倍数与增益下降倍数相同}
\end{equation}

\subsection{减小非线性失真}

非线性失真是由器件特性的非线性产生的。引入负反馈后，若输出出现失真，反馈量中也含有同样的失真成分，它与输入相减使净输入量产生"预失真"，从而抵消一部分输出失真。定量结论：在\textbf{信号源潜力足够}（能提供闭环后所需的净输入量）且失真\textbf{产生于反馈环内部}的条件下，非线性失真减小到 $1/(1+AF)$。
\begin{itemize}
    \item 若失真来自\textbf{外界干扰}（如电源纹波窜入输出端），负反馈无法抑制；
    \item 若信号源内阻过大、无法提供足够的净输入量，则失真不能被有效抑制。
\end{itemize}

\subsection{改变输入电阻与输出电阻}

串联负反馈使输入电阻增大，并联负反馈使输入电阻减小：
\begin{equation}
R_{if}=(1+AF)R_i\quad(\text{串联}),\qquad
R_{if}=\frac{R_i}{1+AF}\quad(\text{并联})
\end{equation}
电压负反馈稳定输出电压，故使输出电阻减小；电流负反馈稳定输出电流，故使输出电阻增大：
\begin{equation}
R_{of}=\frac{R_o}{1+AF}\quad(\text{电压}),\qquad
R_{of}=(1+AF)R_o\quad(\text{电流})
\end{equation}
上述结论都有明确前提：$R_{if}$ 指反馈环内的输入电阻，若反馈网络本身并在输入端，还要与反馈网络的电阻合并计算；$R_{of}$ 指反馈环内的输出电阻。

\subsection{引入负反馈的一般原则}

\begin{enumerate}
    \item 要\textbf{稳定静态工作点}，应引入直流负反馈；
    \item 要\textbf{改善动态性能}（稳定增益、展宽频带、减小失真），应引入交流负反馈；
    \item 要\textbf{提高输入电阻}用串联负反馈，要\textbf{降低输入电阻}用并联负反馈；
    \item 要\textbf{稳定输出电压}（减小输出电阻、增强带负载能力）用电压负反馈，要\textbf{稳定输出电流}（增大输出电阻、作恒流源）用电流负反馈；
    \item 信号源内阻大时宜用串联负反馈（否则反馈电流被信号源内阻分流而削弱），内阻小时宜用并联负反馈。
\end{enumerate}

\subsection{负反馈与自激振荡}

由 $A_f=A/(1+AF)$，当 $1+AF=0$ 时电路自激。由于 $A$、$F$ 都是频率的函数：低频与中频段负反馈的附加相移为零，$1+AF>1$；在高频段或低频段每级放大电路最多可贡献 $90^\circ$ 的附加相移，多级电路累计附加相移达到 $180^\circ$ 时负反馈就变成了正反馈，若此时幅值条件也满足便产生自激。

\begin{equation}
1+\dot A\dot F=0
\ \Longleftrightarrow\
\begin{cases}
|\dot A\dot F|=1 & \text{（幅值条件）}\\
\varphi_A+\varphi_F=\pm180^\circ & \text{（相位条件）}
\end{cases}
\formulaNote{两者同时满足才振荡}
\end{equation}
工程上用波特图判断稳定性：在 $20\lg|\dot AF|=0\,\text{dB}$ 的频率处，若附加相移的绝对值小于 $180^\circ$，则电路稳定；两者的差额称为\textbf{相位裕度}，一般要求不小于 $45^\circ$。若电路已自激，可减小反馈深度或在反馈网络中加入补偿电容（校正）来破坏幅值条件或相位条件。

\paragraph{本章小结}
\begin{itemize}
    \item 四个判断的顺序是：有无反馈 $\to$ 正/负（瞬时极性法） $\to$ 直流/交流（看反馈通路有无电容） $\to$ 组态（输出短路法 + 输入端比较方式）；
    \item 一切定量结论都源自 $A_f=A/(1+AF)$：深度负反馈下 $A_f\approx1/F$，性能改善倍数都是 $(1+AF)$；
    \item 负反馈不能提高增益、不能抑制外界干扰、不能消除由输入信号本身携带的失真，这三点是最常见的概念性错误。
\end{itemize}
